\documentclass[10pt, letterpaper]{article}

% ==== Packages & setup ====
\usepackage[
    ignoreheadfoot,
    top=0.5 cm,
    bottom=0.5 cm,
    left=0.5 cm,
    right=0.5 cm,
    footskip=1.0 cm
]{geometry}

\usepackage[explicit]{titlesec}
\usepackage{tabularx}
\usepackage{array}
\usepackage[dvipsnames]{xcolor}
\definecolor{primaryColor}{RGB}{0, 0, 0}
\usepackage{enumitem}
\usepackage{fontawesome5}
\usepackage{amsmath}
\usepackage[hidelinks]{hyperref}

\usepackage[pscoord]{eso-pic}
\usepackage{calc}
\usepackage{bookmark}
\usepackage{lastpage}
\usepackage{changepage}
\usepackage{paracol}
\usepackage{ifthen}
\usepackage{needspace}
\usepackage{iftex}

% Ensure ATS/machine-readable
\ifPDFTeX
    \input{glyphtounicode}
    \pdfgentounicode=1
    \usepackage[T1]{fontenc}
    \usepackage[utf8]{inputenc}
    \usepackage{lmodern}
\fi

\usepackage[default, type1]{sourcesanspro}
\usepackage{amssymb}

% ==== Global tightness tweaks (keeps it to one page) ====
\linespread{0.97}
\setlength{\tabcolsep}{4pt}
\setlength{\parskip}{0pt}
\AtBeginEnvironment{adjustwidth}{\partopsep0pt}
\pagestyle{empty}
\setcounter{secnumdepth}{0}
\setlength{\parindent}{0pt}
\setlength{\topskip}{0pt}
\setlength{\columnsep}{0.10cm}
\makeatletter

% Section title: smaller size, tighter paddings, thinner rule
\titleformat{\section}{
    \needspace{2\baselineskip}
    \large\color{primaryColor}
}{}{}{
    \textbf{#1}\hspace{0.12cm}\titlerule[0.6pt]\hspace{-0.08cm}
}[]
\titlespacing{\section}{-1pt}{0.15 cm}{0.12 cm}

% Tight bullets
\newenvironment{highlights}{
    \begin{itemize}[
        topsep=0pt,
        parsep=0pt,
        partopsep=0pt,
        itemsep=0pt,
        leftmargin=0.30 cm + 8pt
    ]
}{
    \end{itemize}
}
\newenvironment{highlightsforbulletentries}{
    \begin{itemize}[
        topsep=0pt,
        parsep=0pt,
        partopsep=0pt,
        itemsep=0pt,
        leftmargin=8pt
    ]
}{
    \end{itemize}
}

% Tight single/dual column blocks
\newenvironment{onecolentry}{
    \begin{adjustwidth}{0.18 cm}{0.18 cm}
}{
    \end{adjustwidth}
}
\newenvironment{twocolentry}[2][]{
    \onecolentry
    \def\secondColumn{#2}
    \setcolumnwidth{\fill, 4.2 cm}
    \begin{paracol}{2}
}{
    \switchcolumn \raggedleft \secondColumn
    \end{paracol}
    \endonecolentry
}
\newenvironment{threecolentry}[3][]{
    \onecolentry
    \def\thirdColumn{#3}
    \setcolumnwidth{0.9 cm, \fill, 4.2 cm}
    \begin{paracol}{3}
    {\raggedright #2} \switchcolumn
}{
    \switchcolumn \raggedleft \thirdColumn
    \end{paracol}
    \endonecolentry
}

% Header
\newenvironment{header}{
    \setlength{\topsep}{-0.15cm}\par\kern\topsep\centering\linespread{1.16}
}{
    \par\kern\topsep
}

% Restore normal \href
\let\hrefWithoutArrow\href
\renewcommand{\href}[2]{\hrefWithoutArrow{#1}{#2}}

\definecolor{gr}{RGB}{225,225,225}

\begin{document}

\newcommand{\AND}{\unskip
    \cleaders\copy\ANDbox\hskip\wd\ANDbox
    \ignorespaces
}
\newsavebox\ANDbox
\sbox\ANDbox{}

% ===== Header =====
\begin{header}
    \fontsize{24}{26}\selectfont\scshape Connor Savugot
\end{header}
\begin{center}
{\fontsize{9}{11}\selectfont
\href{https://maps.app.goo.gl/6EoyjWCwyq2M3qoJ6}{\faMapMarker*\enspace Murphy, NC}\hspace{0.9em}%
\href{mailto:consav04@gmail.com}{\faEnvelope\enspace consav04@gmail.com}\hspace{0.9em}%
\href{tel:+1-828-541-9487}{\faPhone\enspace +1 (828) 541-9487}\hspace{0.9em}%
\href{https://linkedin.com/in/connorsavugot}{\faLinkedinIn\enspace connorsavugot}\hspace{0.9em}%
\href{https://github.com/Clem085}{\faGithub\enspace Clem085}
}
\end{center}

% ===== Objective =====
\section{Objective}
\begin{onecolentry}
Seeking an entry-level role in embedded systems programming, firmware engineering, or software engineering. \newline Focus: low-level C/C++, RTOS, and firmware/BIOS. Available May 2026; open to relocation.
\end{onecolentry}
\vspace{7pt}

% ===== Education =====
\section{EDUCATION}
\begin{onecolentry}
\begin{tabularx}{\textwidth}{@{}X r@{}}
    \textbf{North Carolina State University} & \textit{May 2026 \;\;|\;\; Raleigh, NC} \\
    \textit{Bachelor of Science in Computer Engineering} & \textbf{\textit{GPA: 3.0}} \\
\end{tabularx}
\end{onecolentry}
\vspace{5pt}

% ===== Leadership (IEEE) =====
\begin{twocolentry}{Fall 2025 -- Present}
    \textbf{IEEE at NC State \;| \; ECE 306 Parts Distribution Coordinator \& OPS2 Instructor}
    \begin{highlights}
        \item[\textbf{$\vartriangleright$}] Leading the ECE 306 parts program by centralizing sourcing (bulk buys, vetted substitutions), consolidating shipping, and running on-campus distribution. Raised \$2{,}000 for the IEEE chapter while lowering student costs.
        \item[\textbf{$\vartriangleright$}] Teaching OPS2 -- Open Project Space 2 by guiding 10-week student projects, teaching through-hole and SMD soldering, embedded-systems programming, configuring PlatformIO in VS Code, and using GitHub for version control.
    \end{highlights}
\end{twocolentry}


\vspace{7pt}

% ===== Skills =====
\section{SKILLS}
\begin{onecolentry}
\textbf{Programming:} C, C++, Python, Verilog, Bash, MATLAB, VBA, LC-3 Assembly \\
\textbf{Embedded Systems:} TI TMS320 DSP, PIC32, MSP430, PlatformIO, RTOS concepts \\
\textbf{Tools:} Git/GitHub (GitFlow, release tagging), VS Code, Linux CLI \\
\textbf{Hardware:} Soldering (through-hole \& SMD), PCB design basics, microcontrollers, circuit analysis
\end{onecolentry}


\vspace{7pt}
% ===== Work Experience =====
\section{WORK EXPERIENCE}

\begin{twocolentry}{Summers 2024 \& 2025}
    \textbf{Aegis Power Systems \;| \; Embedded Systems Firmware Intern}
    \begin{highlights}
        \item[\textbf{$\vartriangleright$}] Developed PWM firmware in C on a TI TMS320 DSP. Configured 12 channels with per-channel frequency (kHz), duty cycle, and phase shift. Enabled A/B mirroring or inversion and added an ADC-read potentiometer to live-adjust PWM6 duty cycle.
        \item[\textbf{$\vartriangleright$}] Built a remote CLI over SSH for embedded power supplies and integrated RS232 and CAN/J1939 for real-time control.
        \item[\textbf{$\vartriangleright$}] Designed a two-processor update path where the PIC32 managed I2C, SPI, CAN/J1939, and UART orchestration while the TMS320 DSP executed PWM control.
        \item[\textbf{$\vartriangleright$}] Wrote a C++ serial flasher for the DSP and a Python TCP/IP sender for the PIC32 with chunking and checksum verification to prevent corrupted updates.
        \item[\textbf{$\vartriangleright$}] Implemented a fail-safe updater. The PIC32 stored images in SPI flash, maintained a golden image, reflashed the DSP over UART, and committed verified firmware to on-chip flash.
        \item[\textbf{$\vartriangleright$}] Standardized Git practices (GitFlow, signed annotated tags) and authored an 80+ page illustrated guide, handout, and training slides.
    \end{highlights}
\end{twocolentry}
\vspace{5pt}
\begin{twocolentry}{Summers 2021 \& 2022}
    \textbf{TEAM Industries \;| \; Engineering Intern}
    \begin{highlights}
        \item[\textbf{$\vartriangleright$}] Built a VBA quality-check tool in Excel that ingested measurement data, compared it against tolerance tables, and flagged out-of-spec parts with instant summaries.
        \item[\textbf{$\vartriangleright$}] Engineered robust spreadsheets with custom formulas, data validation, and pivot summaries that reduced manual errors and made updates repeatable.
        \item[\textbf{$\vartriangleright$}] Authored clear change reports with before-and-after specs, rationale, and action items that improved revision traceability across engineering and suppliers.
    \end{highlights}
\end{twocolentry}

\vspace{7pt}
% ===== Projects =====
\section{Projects}
\begin{twocolentry}{Fall 2024}
    \textbf{Embedded Systems Toy Car (MSP430, C, CCS, PCB)}
    \begin{highlights}
        \item[\textbf{$\vartriangleright$}] Built and programmed a toy car using a TI MSP430 microcontroller and custom PCBs; integrated IoT Wi-Fi, ADC line following, and PWM motor control.
        \item[\textbf{$\vartriangleright$}] Used an LCD for real-time diagnostics and two programmable buttons plus a rotary dial to select inputs and switch operation modes.

    \end{highlights}
\end{twocolentry}

\vspace{5pt}

\begin{twocolentry}{Summer 2024}
    \textbf{Linux Multi-Boot \& GPU Power Management (Personal Project)}
    \begin{highlights}
        \item[\textbf{$\vartriangleright$}] Partitioned a laptop with Garuda (Arch-based), Ubuntu, and Windows; configured GRUB to detect all systems and set Garuda as default.
        \item[\textbf{$\vartriangleright$}] Wrote \texttt{systemd} service scripts to detect HDMI connections and toggle iGPU/dGPU automatically, extending battery life when an external display was disconnected.
    \end{highlights}
\end{twocolentry}

\vspace{5pt}

\begin{twocolentry}{Spring 2024}
    \textbf{C++ Maze Solver \& Search Algorithms}
    \begin{highlights}
        \item[\textbf{$\vartriangleright$}] Implemented BFS, DFS, and Greedy Best-First search; benchmarked runtime and memory tradeoffs; documented results against classical baselines.
    \end{highlights}
\end{twocolentry}

\end{document}
